\section{Ordered representation and coordinate mapping}
\label{sec:gauge}

\subsection{Linear-time decoding}

For $v_i\leq i$, initialize a component-root array by $r[x]=x$.
Process $i=N-2,\ldots,0$.
At step $i$, join the components with roots $r[v_i]$ and $r[i+1]$ at a new node
$m=2N-2-i$, then replace $r[v_i]$ by $m$.
Before this step, the active components have distinct minimum labels $0,\ldots,i+1$.
Joining the two specified components removes minimum label $i+1$ and preserves minimum label $v_i$.
Induction therefore establishes the invariant and the correctness of the construction.
Each step uses constant work, giving an $O(N)$ decoder.
This argument depends on the ordered condition and does not apply unchanged to a general vector with $v_i>i$.

\subsection{Relabeling a decoded tree}

\begin{proposition}[Ordered representation after relabeling]
\label{thm:relabel}
Given a decoded rooted binary tree, a parent-before-child ordering of its internal nodes constructs an ordered vector and a leaf permutation in $O(N)$ time.
The represented labeled tree and its BME score are unchanged.
\end{proposition}
\begin{proof}
Give the root temporary label zero and initialize a counter $t=1$.
At each internal node, choose an ordering $(a,b)$ of its children.
Pass the node's label to $a$, give $b$ label $t$, record $v_{t-1}$ as the inherited label, and increment $t$.
Each newly issued label follows a unique chain of inheriting children to a leaf.
Every leaf receives exactly one label, so the resulting labels form a permutation of $0,\ldots,N-1$.

The inherited label was issued at an ancestor and is smaller than $t$, hence $v_{t-1}\leq t-1$.
To verify decoding, process the recorded splits in reverse order.
All descendants of a node were assigned their new labels later than that node, so both child components have already been assembled when its recorded merge is reached.
The merge therefore recreates that node's original clade.
Induction reconstructs the entire tree.
Permuting both axes of $D$ by the inverse leaf-label map preserves all terms in Equation~\eqref{eq:pauplin}.
Each vertex is processed a constant number of times.
\end{proof}

The implementation chooses the child order from the minimum original leaf label in each child component, with a second orientation reversing that choice.
These are two deterministic coordinate systems for the same current tree.
Their existence proves representational coverage after suitable relabeling, not coverage by a particular optimizer trajectory.
For general candidate vectors, the implementation also uses a tree-to-vector encoder.
Its unrooted version first inserts a degree-two root on the edge incident to leaf zero, so round-trip correctness is assessed by unrooted splits rather than root placement.

\subsection{Mapping coordinate values to regraft edges}
\label{sec:mapping}

The ordered decoder numbers leaves $0,\ldots,N-1$ and internal nodes $N,\ldots,2N-2$.
The last index is the artificial root.
Suppressing it leaves a computational root $R=2N-3$ and the top edge joining $R$ to the artificial root's other child, denoted $c_{\mathrm{other}}$.
Index ordinary edges by their child endpoints.

For row $i$, let $E_i[u]$ be the score change from pruning the component associated with leaf $i+1$ and regrafting it on the edge indexed by $u$.
The half-tree SPR traversal supplies these values.
Let $b_i(j)$ be the root of the component containing leaf $j$ after the merges with coordinate index greater than $i$.
Then the edge-to-coordinate map is
\begin{equation}
 \Delta_{i,j}=
 \begin{cases}
 E_i[b_i(j)], &0\leq j\leq i,\\[2pt]
 E_i[(2N-2)-2i+j], &i+1\leq j\leq2i-2,\\[2pt]
 E_i[c_{\mathrm{other}}], &j\in\{2i-1,2i\}\cap[i+1,2i].
 \end{cases}
 \label{eq:full_mapping}
\end{equation}
The first case identifies the component on the corresponding pendant edge of the prefix tree.
In the ordered decoder's node numbering, the remaining non-root prefix edges occupy consecutive indices $2N-1-i,\ldots,2N-4$, giving the affine map in the second case.
The last two rooted edge choices collapse onto the same unrooted top edge.
At $i=1$, the middle case is empty and $j=1$ belongs to the first case.
Only $j=2$ takes the collapsed-root value.
Coordinate zero has no nontrivial alternative.

On a CPU, the array $b_i$ can be maintained incrementally as the merge order advances.
On a GPU, independent rows reconstruct their mappings using the latest eligible merge for each minimum leaf.
When several merges have the same minimum leaf, only the latest eligible merge may write that destination.
This avoids concurrent writes that would otherwise depend on thread scheduling.

\section{Parallel dependencies and workspace}
\label{sec:parallel_details}

\subsection{Column ownership and ancestor traversal}

For a fixed column $w$, the bottom-up update in Equation~\eqref{eq:bottom_up} reads only earlier entries of that column.
Thus the recurrence has no cross-column dependency once its initial entries are available.
The implementation separates the necessary transpose from the recurrence rather than writing both triangular copies concurrently.
Top-down correction visits ancestors of $w$ in order from the computational root.
If a node's children have contiguous depth-first intervals, comparing the entry time of $w$ with their interval boundary selects the child containing $w$.
Only the current node and the fixed target are required as traversal state.
The initial degree-three root update uses the other two root branches.

\subsection{Autonomous row reconstruction}

Store the balanced distances from internal vertices to leaves.
For a requested row $w$, these entries supply its leaf values.
For leaf $w$, the leaf values instead come directly from $D$.
Average child entries bottom-up to fill internal positions, then overwrite ancestors of $w$ top-down using the sibling and parent directions.
Initialize the degree-three root correction explicitly and set the diagonal to zero.

For an internal position $u$ in a disjoint component, the bottom-up average is the required balanced distance by linearity.
For a descendant $u$ of $w$, the stored leaf values already use the component at $w$ directed away from that descendant, so averaging within the descendant gives the correct value.
Ancestors of $w$ require the complementary direction and are supplied by the top-down overwrite.
These cases exhaust the row.
The reconstruction costs $O(N+\operatorname{depth}(w))=O(N)$.
Reconstructing a constant number of rows per pruned component therefore retains quadratic sweep complexity.
The approximately fourfold reduction concerns the balanced-distance table only.

\subsection{Traversal stack size}

When a traversal reaches two outward components, push the larger component first and the smaller second on a last-in-first-out stack.
The smaller component is processed next.
Each extra pending larger component is associated with descent into a component at most half the size of its predecessor.
Hence the number of simultaneously pending components is logarithmic in the number of tree vertices.
Including the active traversal state gives a safe bound
$\lfloor\log_2(2N-2)\rfloor+2$.
This is a workspace bound, not a statement that the traversal visits fewer vertices.

\section{Why profile completeness and current outside information matter}
\label{sec:counterexamples}

Two constructions illustrate the assumptions in Section~\ref{sec:acceptance}.
On a 4,102-leaf tree containing a balanced 4,096-leaf component, an adversarial split distance can give a capped proposal change of $-0.25$ while its true full-profile change is $+0.2421875$ for a cap of 64, or $+0.234375$ for a cap of 128.
The proposal sign alone is therefore insufficient.
Full-profile acceptance rejects both changes.

Separately, consider two disjoint windows whose internal rearrangements change their root-to-leaf weight vectors from $p_A,p_B$ to $p_A+\delta p_A,p_B+\delta p_B$.
With a fixed path coefficient $\kappa$ between their attachment points, their cross contribution contains
$\kappa p_A^\top D_{AB}p_B$.
The jointly changed contribution includes the additional term
\[
 \kappa\,\delta p_A^\top D_{AB}\delta p_B,
\]
which is absent from the sum of two changes computed against the original surroundings.
In a 1,540-leaf test with two 768-leaf windows, changes
$-0.0032945313$ and $-0.0254489835$ computed against frozen surroundings combine to an actual increase of $+0.0505258348$.
Rebuilding the outside profile after the first commit makes the other proposal positive:
$+0.0538203660$ or $+0.0759748183$, depending on the processing order.
It is then rejected.
These values come from controlled CPU replay of the acceptance and window-controller paths.
They do not constitute a native-GPU concurrency test.

\section{Implementation and reproducibility}
\label{sec:software}

The CPU evaluator, tree encoders, and full-tree optimizer are implemented in
\texttt{phylo2vec\_fastme.c}, with input validation and Python bindings in
\texttt{discrete\_gradient.py}.
The lower-triangular and full-coordinate interfaces require an ordered baseline.
The multi-GPU evaluator and its tiled variants are implemented in
\texttt{phylo2vec\_cuda\_8xh100.cu}.
Sequence refinement is exercised by \texttt{bench\_rnasim\_gpu.cu}.
The distinct distance-oracle recovery workload is implemented in
\texttt{bench\_minibatch\_recovery.cu}.

Independent checks include \texttt{verify\_all.c} for objective identities and encoding boundaries,
\texttt{verify\_spr\_coverage.py} for neighborhood enumeration, and
\texttt{compare\_with\_fastme\_cli.py} for the external reference.
The profile-completeness and window-composition constructions are recorded with their reproduction commands in the discrete-gradient review artifacts.
The fifth review checks the window controller, full-profile arithmetic, and a shared-input batching path using CPU emulation.
That harness does not cover the subsequent split-launch implementation.
Native CUDA tests and the implementation revision associated with each timing must be recorded separately.

The present manuscript reorganizes existing measurements rather than reporting a new benchmark campaign.
The source archive and revision notes accompanying this draft preserve the original tables and identify the corrected sequence timings used in the main comparison.
Before submission, benchmark records should be consolidated with source revisions, compiler flags, seeds, repeat counts, and the exact timing boundaries.
The earlier recovery and controlled-JC69 minibatch timings are retained in that archive rather than used to characterize the corrected acceptance path here.

\clearpage
\section{Supplementary CPU comparisons}
\label{sec:cpu_tables}

Tables~\ref{tab:bench_noiseless} and~\ref{tab:bench_noisy} give the full Apple M5 comparisons.
FME denotes the in-memory reference implementation.
P2V denotes Phylo2Vec.
RF denotes normalized Robinson--Foulds distance.
Wins and ties refer to the BME objective, with tolerance $10^{-9}$.
Timings and evaluation counts are means, with standard deviations and the reported P2V median and interquartile range retained.

\begin{table}[!htbp]
\centering\small
\setlength{\tabcolsep}{4pt}
\caption{Additive distance comparison on Apple M5.}
\label{tab:bench_noiseless}
\par\smallskip
\textit{Objective and topology}\par\smallskip
\begin{tabular}{@{}rrrrrr@{}}
\toprule
$N$ & FME $L$ & P2V $L$ & Win/tie/loss & FME RF & P2V RF\\
\midrule
$20$ & $11.12739$ & $11.12739$ & \textbf{0 / 20 / 0} & $0.0000$ & $0.0000$\\
$50$ & $29.42572$ & $29.42572$ & \textbf{0 / 20 / 0} & $0.0000$ & $0.0000$\\
$100$ & $59.52693$ & $59.52693$ & \textbf{0 / 10 / 0} & $0.0000$ & $0.0000$\\
$150$ & $89.04365$ & $89.04365$ & \textbf{0 / 5 / 0} & $0.0000$ & $0.0000$\\
$200$ & $116.32398$ & $116.32398$ & \textbf{0 / 5 / 0} & $0.0000$ & $0.0000$\\
\bottomrule
\end{tabular}
\par\smallskip
\textit{Computational cost}\par\smallskip
\begin{tabular}{@{}rrrrrrr@{}}
\toprule
$N$ & FME iters. & P2V evals. & FME ms & P2V ms & P2V med. [IQR] & Ratio\\
\midrule
$20$ & $13.4$ & $11.2$ & $0.12 \pm 0.02$ & $0.08 \pm 0.01$ & $0.07\,[0.02]$ & \textbf{1.54$\times$}\\
$50$ & $42.2$ & $19.4$ & $1.58 \pm 0.35$ & $0.57 \pm 0.13$ & $0.59\,[0.16]$ & \textbf{2.78$\times$}\\
$100$ & $97.7$ & $35.4$ & $9.05 \pm 0.84$ & $2.63 \pm 0.26$ & $2.65\,[0.44]$ & \textbf{3.44$\times$}\\
$150$ & $156.2$ & $49.4$ & $31.47 \pm 0.77$ & $7.92 \pm 0.86$ & $7.75\,[0.64]$ & \textbf{3.97$\times$}\\
$200$ & $211.0$ & $63.6$ & $75.46 \pm 2.42$ & $18.14 \pm 0.74$ & $18.05\,[1.17]$ & \textbf{4.16$\times$}\\
\bottomrule
\end{tabular}

\end{table}

\begin{table}[!htbp]
\centering\small
\setlength{\tabcolsep}{4pt}
\caption{Noisy distance comparison on Apple M5.}
\label{tab:bench_noisy}
\par\smallskip
\textit{Objective and topology}\par\smallskip
\begin{tabular}{@{}rrrrrr@{}}
\toprule
$N$ & FME $L$ & P2V $L$ & Win/tie/loss & FME RF & P2V RF\\
\midrule
$20$ & $11.19777$ & $11.19777$ & \textbf{0 / 20 / 0} & $0.0500$ & $0.0500$\\
$50$ & $29.25507$ & $29.25607$ & \textbf{1 / 18 / 1} & $0.0660$ & $0.0649$\\
$100$ & $59.25471$ & $59.25471$ & \textbf{0 / 10 / 0} & $0.0526$ & $0.0526$\\
$150$ & $92.14846$ & $92.14846$ & \textbf{0 / 5 / 0} & $0.0395$ & $0.0395$\\
$200$ & $120.94220$ & $120.94220$ & \textbf{1 / 4 / 0} & $0.0467$ & $0.0487$\\
\bottomrule
\end{tabular}
\par\smallskip
\textit{Computational cost}\par\smallskip
\begin{tabular}{@{}rrrrrrr@{}}
\toprule
$N$ & FME iters. & P2V evals. & FME ms & P2V ms & P2V med. [IQR] & Ratio\\
\midrule
$20$ & $14.4$ & $12.4$ & $0.07 \pm 0.01$ & $0.05 \pm 0.01$ & $0.04\,[0.02]$ & \textbf{1.60$\times$}\\
$50$ & $46.1$ & $25.9$ & $1.07 \pm 0.07$ & $0.49 \pm 0.09$ & $0.49\,[0.10]$ & \textbf{2.18$\times$}\\
$100$ & $102.6$ & $41.0$ & $10.62 \pm 2.53$ & $3.32 \pm 0.56$ & $3.18\,[0.52]$ & \textbf{3.19$\times$}\\
$150$ & $166.0$ & $61.0$ & $33.97 \pm 1.01$ & $10.15 \pm 1.35$ & $10.17\,[1.35]$ & \textbf{3.35$\times$}\\
$200$ & $235.8$ & $84.4$ & $85.16 \pm 2.06$ & $24.67 \pm 2.79$ & $24.51\,[4.48]$ & \textbf{3.45$\times$}\\
\bottomrule
\end{tabular}

\end{table}


\clearpage
\section{Supplementary scaling measurements}
\label{sec:scaling_tables}

The tables below retain the reported sweep measurements, separating CPU/Metal, CUDA, sharded, and embedding-based workloads.
The single-GPU speedups use their respective host CPU references.
They should not be interpreted as comparisons against one common host baseline.
Stage timings measure balanced-distance construction and SPR evaluation, rather than optimization to convergence.
The $2N^2$ column in the sharded table is the original rectangular-workload size proxy.
It is not the number of valid Phylo2Vec coordinate entries, which is $(N-1)^2$, nor a count of distinct neighboring trees.
Reported memory units are retained from the benchmark records.

For the embedding workload, each batch has 20,000 skeleton leaves, comprising 15,000 expanded leaves and 5,000 condensed components.
Its epoch columns multiply average condensation-plus-gradient time by
$\lceil N/15{,}000\rceil$.
They are workload estimates rather than measured convergence times.
The largest listed validation discrepancy in this workload is $1.4\times10^{-5}$, and the listed validation outcomes refer to the benchmark's own checks.

\begin{table}[!htbp]
\centering\small
\setlength{\tabcolsep}{4pt}
\caption{Single-GPU gradient timings and agreement with the CPU reference. Times are milliseconds.}
\label{tab:cuda_a6000_verify}
\par\smallskip
\textit{RTX A6000 timings and reported maximum errors}\par\smallskip
\begin{tabular}{@{}rrrrr@{}}
\toprule
$N$ & A6000 full & Ratio & A6000 tiled & Max. error\\
\midrule
$100$ & $0.061$ & \textbf{8.44$\times$} & $0.078$ & $3.68\text{e-}07$\\
$250$ & $0.186$ & \textbf{15.92$\times$} & $0.207$ & $6.22\text{e-}07$\\
$500$ & $0.382$ & \textbf{32.95$\times$} & $0.511$ & $6.48\text{e-}07$\\
$1{,}000$ & $1.035$ & \textbf{33.69$\times$} & $1.988$ & $9.29\text{e-}07$\\
$2{,}000$ & $3.851$ & \textbf{40.15$\times$} & $7.288$ & $1.11\text{e-}06$\\
$4{,}000$ & $11.053$ & \textbf{62.48$\times$} & $43.124$ & $1.07\text{e-}06$\\
\bottomrule
\end{tabular}
\par\smallskip
\textit{H100 timings}\par\smallskip
\begin{tabular}{@{}rrrrrrr@{}}
\toprule
$N$ & CPU ref. & H100 avg. & H100 SPR & H100 total & Ratio & H100 tiled\\
\midrule
$100$ & $0.249$ & $0.044$ & $0.027$ & $0.071$ & \textbf{3.51$\times$} & $0.095$\\
$250$ & $1.964$ & $0.116$ & $0.048$ & $0.164$ & \textbf{11.98$\times$} & $0.205$\\
$500$ & $8.176$ & $0.227$ & $0.077$ & $0.304$ & \textbf{26.89$\times$} & $0.460$\\
$1{,}000$ & $32.336$ & $0.661$ & $0.105$ & $0.766$ & \textbf{42.21$\times$} & $1.030$\\
$2{,}000$ & $154.091$ & $1.557$ & $0.726$ & $2.283$ & \textbf{67.50$\times$} & $4.273$\\
$4{,}000$ & $701.375$ & $3.565$ & $2.654$ & $6.219$ & \textbf{112.78$\times$} & $22.327$\\
\bottomrule
\end{tabular}
\par\medskip\begin{minipage}{\linewidth}\footnotesize The maximum-error column is the reported combined check across the GPU paths. All listed checks pass their benchmark tolerance. CPU baselines differ between hosts.\end{minipage}
\end{table}

\begin{table}[!htbp]
\centering\small
\setlength{\tabcolsep}{4pt}
\caption{Embedding-based condensed sweeps with $K=20{,}000$ skeleton leaves. The epoch columns are batch-count times mean condensation-plus-gradient time.}
\label{tab:minibatch_scaling}
\par\smallskip
\textit{Two A6000 GPUs}\par\smallskip
\begin{tabular}{@{}rrrr@{}}
\toprule
$N$ & Batches & 2 A6000 cond.+grad. (ms) & Epoch (s)\\
\midrule
$1{,}000{,}000$ & $67$ & $4.45 + 284.44$ & $19.36$\\
$5{,}000{,}000$ & $334$ & $8.42 + 283.02$ & $97.34$\\
$10{,}000{,}000$ & $667$ & $17.29 + 281.58$ & $199.35$\\
\bottomrule
\end{tabular}
\par\smallskip
\textit{Eight H100 GPUs}\par\smallskip
\begin{tabular}{@{}rrrrrr@{}}
\toprule
$N$ & Cond. ms & Grad. ms & Epoch s & MB/GPU & Check\\
\midrule
$1{,}000{,}000$ & $2.69$ & $48.53$ & $\mathbf{3.43}$ & $7{,}793$ & \textbf{PASS} ($1.4\text{e-}05$)\\
$5{,}000{,}000$ & $4.86$ & $48.48$ & $\mathbf{17.82}$ & $7{,}824$ & \textbf{PASS} ($2.1\text{e-}06$)\\
$10{,}000{,}000$ & $10.20$ & $48.57$ & $\mathbf{39.20}$ & $7{,}862$ & \textbf{PASS} ($2.2\text{e-}06$)\\
\bottomrule
\end{tabular}
\par\medskip\begin{minipage}{\linewidth}\footnotesize This workload uses a bilinear embedding distance and is separate from sequence-distance refinement. The validation figures are observed discrepancies, not universal error bounds.\end{minipage}
\end{table}

\begin{table}[!htbp]
\centering\small
\setlength{\tabcolsep}{4pt}
\caption{Full-table and tiled multi-GPU sweeps. Times are milliseconds and memory is reported in MB.}
\label{tab:cuda_a6000_scaling}
\par\smallskip
\textit{A6000 timings and memory columns}\par\smallskip
\begin{tabular}{@{}rrrrr@{}}
\toprule
$N$ & 2 A6000 full & 2 A6000 tiled & Full memory & Tiled memory\\
\midrule
$1{,}000$ & $0.93$ & $5.75$ & $27$ & $39$\\
$2{,}000$ & $4.89$ & $13.90$ & $92$ & $93$\\
$4{,}000$ & $14.69$ & $56.85$ & $336$ & $247$\\
$10{,}000$ & $90.45$ & $249.59$ & $1{,}985$ & $1{,}076$\\
$20{,}000$ & $422.73$ & $1{,}212.27$ & $7{,}785$ & $3{,}677$\\
$40{,}000$ & \emph{OOM} & $5{,}628.62$ & $31{,}142$ & $13{,}614$\\
\bottomrule
\end{tabular}
\par\smallskip
\textit{Full-table H100 sweep}\par\smallskip
\begin{tabular}{@{}rrrr@{}}
\toprule
$N$ & 1 H100 [avg.+SPR] & 8 H100 & Ratio\\
\midrule
$1{,}000$ & $0.76\; [0.65 + 0.11]$ & $0.79$ & $0.96\times$\\
$2{,}000$ & $2.21\; [1.47 + 0.74]$ & $1.71$ & \textbf{1.29$\times$}\\
$4{,}000$ & $6.01\; [3.57 + 2.44]$ & $3.98$ & \textbf{1.51$\times$}\\
$10{,}000$ & $29.61\; [10.10 + 19.50]$ & $13.03$ & \textbf{2.27$\times$}\\
$20{,}000$ & $155.93\; [23.23 + 132.70]$ & $41.06$ & \textbf{3.80$\times$}\\
$40{,}000$ & $816.34\; [60.08 + 756.26]$ & $\mathbf{158.36}$ & \textbf{5.15$\times$}\\
\bottomrule
\end{tabular}
\par\smallskip
\textit{Tiled H100 sweep}\par\smallskip
\begin{tabular}{@{}rrrr@{}}
\toprule
$N$ & 1 H100 tiled & 8 H100 tiled & Top-1 check\\
\midrule
$1{,}000$ & $1.46$ & $4.47$ & \textbf{PASS} ($0.00\text{e+}00$)\\
$2{,}000$ & $6.14$ & $5.21$ & \textbf{PASS} ($2.38\text{e-}07$)\\
$4{,}000$ & $30.78$ & $10.98$ & \textbf{PASS} ($0.00\text{e+}00$)\\
$10{,}000$ & $250.69$ & $57.53$ & \textbf{PASS} ($0.00\text{e+}00$)\\
$20{,}000$ & $1{,}123.43$ & $244.13$ & \textbf{PASS} ($4.77\text{e-}07$)\\
$40{,}000$ & $4{,}768.31$ & $\mathbf{1{,}032.83}$ & \textbf{PASS} ($0.00\text{e+}00$)\\
\bottomrule
\end{tabular}

\end{table}

\begin{table}[!htbp]
\centering\small
\setlength{\tabcolsep}{4pt}
\caption{Distributed sweeps on eight H100 GPUs.}
\label{tab:h100_scaling}
\par\smallskip
\textit{Timing}\par\smallskip
\begin{tabular}{@{}rrrrr@{}}
\toprule
$N$ & Work proxy $2N^2$ & Distance ms & SPR ms & Total s\\
\midrule
$40{,}000$ & $3.20 \times 10^9$ & $514.54$ & $517.92$ & $\mathbf{1.032}$\\
$80{,}000$ & $1.28 \times 10^{10}$ & $2{,}112.82$ & $2{,}321.28$ & $\mathbf{4.434}$\\
$120{,}000$ & $2.88 \times 10^{10}$ & $4{,}785.75$ & $4{,}803.60$ & $\mathbf{9.589}$\\
$180{,}000$ & $6.48 \times 10^{10}$ & $10{,}848.67$ & $13{,}977.73$ & $\mathbf{24.826}$\\
$250{,}000$ & $1.25 \times 10^{11}$ & $39{,}981.55$ & $43{,}721.48$ & $\mathbf{83.703}$\\
$350{,}000$ & $2.45 \times 10^{11}$ & $76{,}240.05$ & $96{,}515.83$ & $\mathbf{172.756}$\\
$380{,}000$ & $2.89 \times 10^{11}$ & $89{,}992.00$ & $115{,}457.57$ & $\mathbf{205.450}$\\
\bottomrule
\end{tabular}
\par\smallskip
\textit{Memory and execution mode}\par\smallskip
\begin{tabular}{@{}rrrr@{}}
\toprule
$N$ & MB/GPU & Aggregate GB & Mode and check\\
\midrule
$40{,}000$ & $2{,}386$ & $18.6$ & \texttt{Full-Shard} [\textbf{PASS}]\\
$80{,}000$ & $7{,}824$ & $61.1$ & \texttt{Shard+Calib} [\textbf{PASS}]\\
$120{,}000$ & $16{,}314$ & $127.5$ & \texttt{Shard+Calib} [\textbf{PASS}]\\
$180{,}000$ & $34{,}770$ & $271.6$ & \texttt{Shard+Calib} [\textbf{PASS}]\\
$250{,}000$ & $62{,}296$ & $486.7$ & \texttt{Shard+Calib} [\textbf{PASS}]\\
$350{,}000$ & $62{,}180$ & $485.8$ & \texttt{StreamD+Calib} [\textbf{PASS}]\\
$380{,}000$ & $72{,}946$ & $569.9$ & \texttt{StreamD+Calib} [\textbf{PASS}]\\
\bottomrule
\end{tabular}
\par\medskip\begin{minipage}{\linewidth}\footnotesize At $N=400{,}000$, the reported requirement is 80,600 MB/GPU (629.7 GB aggregate), exceeding available device memory. Full-Shard and Shard+Calib retain distance shards. StreamD+Calib generates distance tiles as required. The work proxy is not a count of distinct moves.\end{minipage}
\end{table}

